\section{Class PSPACE --- TQBFARM --- Walk}
%%% generated by the blueprint pipeline from TCSlib.Complexity.ClassPSPACE.TQBFARM.Walk
\subsection*{Class PSPACE --- TQBFARM --- Walk: Introduction}

Declarations of the Lean module \texttt{TCSlib.Complexity.ClassPSPACE.TQBFARM.Walk}.

\subsection*{Class PSPACE --- TQBFARM --- Walk: Results}

\begin{lemma}[Reachability through an in-range length query]
\lean{Complexity.LogProg.TQBFARM.len_t}
\label{Complexity.LogProg.TQBFARM.len_t}
\leanok
\uses{Complexity.LogProg.AConf, Complexity.LogProg.AReach, Complexity.LogProg.GuardLe, Complexity.LogProg.TQBFARM.A, Complexity.LogProg.TQBFARM.Lb, Complexity.LogProg.TQBFARM.St, Complexity.LogProg.TQBFARM.St.Bd, Complexity.LogProg.TQBFARM.St.toV, Complexity.LogProg.TQBFARM.St.toV_P, Complexity.LogProg.TQBFARM.orc, Complexity.LogProg.TQBFARM.orc_len, Complexity.LogProg.TQBFARM.rP, Complexity.LogProg.TQBFARM.step_of, Complexity.LogProg.astep, Complexity.LogProg.vword_whole₁}
\difficulty{3}
Let $x$ be a finite bit string, $B$ a natural number, $l,l_1,l_0$ machine labels, $st$ a
seven-register state, and $c$ an abstract machine configuration. Suppose the instruction
at $l$ queries whether the whole input length exceeds the pointer register, with $l_1$
the positive-answer label and $l_0$ the negative-answer label. If every register of $st$
is at most $B$, the pointer value in $st$ is less than $|x|$, and there is a run from
label $l_1$ with register values given by $st$ to $c$, with the guard satisfied at every
earlier step, then there is such a guarded run from label $l$ with the same register
values to $c$.
\end{lemma}

\begin{lemma}[Reachability through an out-of-range length query]
\lean{Complexity.LogProg.TQBFARM.len_f}
\label{Complexity.LogProg.TQBFARM.len_f}
\leanok
\uses{Complexity.LogProg.AConf, Complexity.LogProg.AReach, Complexity.LogProg.GuardLe, Complexity.LogProg.TQBFARM.A, Complexity.LogProg.TQBFARM.Lb, Complexity.LogProg.TQBFARM.St, Complexity.LogProg.TQBFARM.St.Bd, Complexity.LogProg.TQBFARM.St.toV, Complexity.LogProg.TQBFARM.St.toV_P, Complexity.LogProg.TQBFARM.orc, Complexity.LogProg.TQBFARM.orc_len, Complexity.LogProg.TQBFARM.rP, Complexity.LogProg.TQBFARM.step_of, Complexity.LogProg.astep, Complexity.LogProg.vword_whole₁}
\difficulty{3}
Let $x$ be a finite bit string, $B$ a natural number, $l,l_1,l_0$ machine labels, $st$ a
seven-register state, and $c$ an abstract machine configuration. Suppose the instruction
at $l$ queries whether the whole input length exceeds the pointer register, with $l_1$
the positive-answer label and $l_0$ the negative-answer label. If every register of $st$
is at most $B$, the pointer value in $st$ is at least $|x|$, and there is a run from
label $l_0$ with register values given by $st$ to $c$, with the guard satisfied at every
earlier step, then there is such a guarded run from label $l$ with the same register
values to $c$.
\end{lemma}

\begin{theorem}[Run of the Syntax Scan]
\lean{Complexity.LogProg.TQBFARM.syn_run}
\label{Complexity.LogProg.TQBFARM.syn_run}
\leanok
\uses{Complexity.LogProg.TQBFARM.A, Complexity.LogProg.TQBFARM.Run, Complexity.LogProg.TQBFARM.St, Complexity.LogProg.TQBFARM.St.Bd, Complexity.LogProg.TQBFARM.SynSt, Complexity.LogProg.TQBFARM.fin_of, Complexity.LogProg.TQBFARM.len_f, Complexity.LogProg.TQBFARM.len_t, Complexity.LogProg.TQBFARM.read_b, Complexity.LogProg.TQBFARM.step_of, Complexity.LogProg.TQBFARM.synStep, Complexity.LogProg.astep, Complexity.getD_of_drop}
\difficulty{5}
Let $x$ be a Boolean word, let $n$ be a natural number, let $q$ be a syntax-automaton
state, and let $st$ be a seven-register state. If $st.P+n=|x|$, every register of $st$
is at most $B$, and $|x|\le B$, then there is a finite run from the syntax-scanning
label with automaton state $q$ and registers $st$ to the input-exhausted label with
automaton state obtained by folding the suffix of $x$ beginning at position $st.P$
through the syntax transition, and with the input pointer set to $|x|$ and every other
register unchanged. Every configuration before the endpoint satisfies the program guard
and register bound $B$.
\end{theorem}

\begin{definition}[State with Prescribed Depth and Pointers]
\lean{Complexity.LogProg.TQBFARM.stp}
\label{Complexity.LogProg.TQBFARM.stp}
\leanok
\uses{Complexity.LogProg.TQBFARM.St}
Given natural numbers $k$, $S$, and $p$, define a state whose depth is $k$, assignment
stack is $S$, and input pointer is $p$, with its scratch copy, copy-loop register,
halving counter, and current-literal variable all equal to zero.
\end{definition}

\begin{lemma}[Invariance under equal register states]
\lean{Complexity.LogProg.TQBFARM.cont}
\label{Complexity.LogProg.TQBFARM.cont}
\leanok
\uses{Complexity.LogProg.AConf, Complexity.LogProg.AReach, Complexity.LogProg.GuardLe, Complexity.LogProg.TQBFARM.A, Complexity.LogProg.TQBFARM.Lb, Complexity.LogProg.TQBFARM.St, Complexity.LogProg.TQBFARM.St.toV, Complexity.LogProg.TQBFARM.orc}
\difficulty{1}
Let $x$ be a finite Boolean word, let $B$ be a natural number, let $st$ and $st'$ be
machine states, let $l$ be a machine label, and let $c$ be an abstract configuration. If
$st=st'$ and the machine reaches $c$ from label $l$ with register vector given by $st'$,
while satisfying the guard that every intermediate configuration has valid preconditions
and all register values at most $B$, then it also reaches $c$ from label $l$ with
register vector given by $st$, under the same guard.
\end{lemma}

\begin{lemma}[Bound on the Initial State]
\lean{Complexity.LogProg.TQBFARM.stp_bd}
\label{Complexity.LogProg.TQBFARM.stp_bd}
\leanok
\uses{Complexity.LogProg.TQBFARM.St.Bd, Complexity.LogProg.TQBFARM.stp}
\difficulty{2}
Let $k,S,p,B$ be natural numbers. If $k\le B$, $S\le B$, and $p\le B$, then the state
with depth $k$, stack value $S$, input pointer $p$, and all four scratch registers equal
to $0$ has every register at most $B$.
\end{lemma}

\begin{lemma}[Bit Positions in an Encoded Pair]
\lean{Complexity.LogProg.TQBFARM.getD_enc}
\label{Complexity.LogProg.TQBFARM.getD_enc}
\leanok
\uses{Complexity.pairEncode_getD, Std.Sat.CNF.serialize, Turing.pairEncode}
\difficulty{1}
If $x=\langle q,\mathrm{serialize}(\varphi)\rangle$, then for every natural number $j$,
the bit at position $j$ of $x$ (taking out-of-range positions to be false) is
$q_{\lfloor j/2\rfloor}$ when $j<2|q|$, is false when $j=2|q|$, is true when $j=2|q|+1$,
and otherwise is the bit of $\mathrm{serialize}(\varphi)$ at position $j-2|q|-2$, again
taking out-of-range positions to be false.
\end{lemma}

\begin{lemma}[Length of the encoded quantified formula]
\lean{Complexity.LogProg.TQBFARM.length_enc}
\label{Complexity.LogProg.TQBFARM.length_enc}
\leanok
\uses{Std.Sat.CNF.serialize, Turing.length_pairEncode, Turing.pairEncode}
\difficulty{2}
Let $x$ and $qs$ be bit strings and let $\varphi$ be a CNF formula over
natural-number-indexed variables. If $x$ is the self-delimiting pair encoding of $qs$
and the serialization of $\varphi$, then $2|qs|+2\leq |x|$.
\end{lemma}

\begin{theorem}[Run through a quantified-formula node]
\lean{Complexity.LogProg.TQBFARM.node_run}
\label{Complexity.LogProg.TQBFARM.node_run}
\leanok
\uses{Complexity.LogProg.AReach.trans, Complexity.LogProg.TQBFARM.A, Complexity.LogProg.TQBFARM.Run, Complexity.LogProg.TQBFARM.cont, Complexity.LogProg.TQBFARM.cpDst, Complexity.LogProg.TQBFARM.cpExit, Complexity.LogProg.TQBFARM.cpSrc, Complexity.LogProg.TQBFARM.cp_run, Complexity.LogProg.TQBFARM.fin_of, Complexity.LogProg.TQBFARM.step_of, Complexity.LogProg.TQBFARM.stp, Complexity.LogProg.TQBFARM.stp_bd, Complexity.LogProg.astep}
\difficulty{4}
For any bit string $x$ and natural numbers $k,S,p,B$ satisfying $2k\le B$, $S\le B$, and
$p\le B$, the seven-register program with its specified oracle has a run from the node
label with depth $k$, stack value $S$, pointer $p$, and all scratch registers zero to
the label $\mathrm{nd1}$ with depth $k$, stack value $S$, pointer $2k$, and all scratch
registers zero. Every configuration before the final configuration satisfies the program
precondition and has every register value at most $B$.
\end{theorem}

\begin{theorem}[Pushing a bit onto the stack]
\lean{Complexity.LogProg.TQBFARM.push_run}
\label{Complexity.LogProg.TQBFARM.push_run}
\leanok
\uses{Complexity.LogProg.AReach.trans, Complexity.LogProg.TQBFARM.A, Complexity.LogProg.TQBFARM.Run, Complexity.LogProg.TQBFARM.St.Bd, Complexity.LogProg.TQBFARM.cont, Complexity.LogProg.TQBFARM.dm_run, Complexity.LogProg.TQBFARM.fin_of, Complexity.LogProg.TQBFARM.mv_run, Complexity.LogProg.TQBFARM.step_of, Complexity.LogProg.TQBFARM.stp, Complexity.LogProg.TQBFARM.stp_bd, Complexity.LogProg.astep}
\difficulty{5}
Let $x$ be a finite bit string, let $B,k,S,p$ be natural numbers, and let $b$ be a
Boolean. If $p\le B$, $2S+1\le B$, and $k+1\le B$, then the seven-register program has a
run, with its specified oracle, from the push state carrying $b$ and register values
$(k,S,p,0,0,0,0)$ to the node state with register values $(k+1,2S+b,0,0,0,0,0)$. Every
configuration strictly before the endpoint satisfies the program precondition and has
all register values at most $B$.
\end{theorem}

\begin{theorem}[Return from a recursive level]
\lean{Complexity.LogProg.TQBFARM.pop_run}
\label{Complexity.LogProg.TQBFARM.pop_run}
\leanok
\uses{Complexity.LogProg.TQBFARM.A, Complexity.LogProg.TQBFARM.Run, Complexity.LogProg.TQBFARM.cont, Complexity.LogProg.TQBFARM.fin_of, Complexity.LogProg.TQBFARM.step_of, Complexity.LogProg.TQBFARM.stp, Complexity.LogProg.TQBFARM.stp_bd, Complexity.LogProg.astep}
\difficulty{5}
For any bit string $x$, Boolean values $v,b$, and natural numbers $k,S,B$ with $k+1\le
B$ and $2S+1\le B$, the seven-register program with its specified oracle has a run from
the return label carrying $v$, at depth $k+1$, stack value $2S+b$, and pointer zero, to
depth $k$, stack value $S$, and pointer zero. The final label is $\mathrm{rt}(v)$ if $b$
is true and the first return-copy label carrying $v$ if $b$ is false. Every
configuration before the final configuration satisfies the program precondition and has
every register value at most $B$.
\end{theorem}

\begin{theorem}[Completing the quantifier lookup]
\lean{Complexity.LogProg.TQBFARM.qlook_run}
\label{Complexity.LogProg.TQBFARM.qlook_run}
\leanok
\uses{Complexity.LogProg.AReach.trans, Complexity.LogProg.TQBFARM.Run, Complexity.LogProg.TQBFARM.cont, Complexity.LogProg.TQBFARM.cpDst, Complexity.LogProg.TQBFARM.cpExit, Complexity.LogProg.TQBFARM.cpSrc, Complexity.LogProg.TQBFARM.cp_run, Complexity.LogProg.TQBFARM.fin_of, Complexity.LogProg.TQBFARM.stp, Complexity.LogProg.TQBFARM.stp_bd}
\difficulty{4}
Let $x$ be a finite bit string, let $B,k,S$ be natural numbers, and let $v$ be a
Boolean. If $2k\le B$ and $S\le B$, then the seven-register program has a run, with its
specified oracle, from the first copy phase of the return procedure carrying $v$ and
register values $(k,S,0,0,0,0,0)$ to the quantifier-reading state carrying $v$ and
register values $(k,S,2k,0,0,0,0)$. Every configuration strictly before the endpoint
satisfies the program precondition and has all register values at most $B$.
\end{theorem}

\begin{definition}[Correctness of a recursive call]
\lean{Complexity.LogProg.TQBFARM.CallOK}
\label{Complexity.LogProg.TQBFARM.CallOK}
\leanok
\uses{Complexity.LogProg.TQBFARM.Run, Complexity.LogProg.TQBFARM.Val, Complexity.LogProg.TQBFARM.stp}
Fix an input bit string $x$, a natural-number bound $B$, a Boolean list $qs$, and a CNF
formula $\varphi$ over natural-number variables. The call is correct at depth $k$ if,
for every assignment stack $S<2^{k+1}$, the machine runs from the node label with
initial state of depth $k$, stack $S$, and pointer $0$ to the return label carrying the
truth value of the quantified formula at $(k,S)$, with the same initial state restored.
\end{definition}

\begin{theorem}[Base case for quantified formula evaluation]
\lean{Complexity.LogProg.TQBFARM.leaf_spec}
\label{Complexity.LogProg.TQBFARM.leaf_spec}
\leanok
\uses{Complexity.LogProg.AReach.trans, Complexity.LogProg.TQBFARM.A, Complexity.LogProg.TQBFARM.CallOK, Complexity.LogProg.TQBFARM.Val, Complexity.LogProg.TQBFARM.asg, Complexity.LogProg.TQBFARM.cont, Complexity.LogProg.TQBFARM.formula_run, Complexity.LogProg.TQBFARM.getD_enc, Complexity.LogProg.TQBFARM.length_enc, Complexity.LogProg.TQBFARM.node_run, Complexity.LogProg.TQBFARM.read_f, Complexity.LogProg.TQBFARM.read_t, Complexity.LogProg.TQBFARM.step_of, Complexity.LogProg.TQBFARM.stp, Complexity.LogProg.TQBFARM.stp_bd, Complexity.LogProg.TQBFARM.val_leaf, Complexity.LogProg.astep, Complexity.LogProg.reach_of_eq, Std.Sat.CNF.serialize, Turing.pairEncode, Turing.pairEncode_eq_dbl}
\difficulty{5}
Let $q$ be a finite Boolean word, let $\varphi$ be a CNF formula over natural-number
variables, and let $x$ be a finite bit string. If $x$ is the pair encoding of $q$ and
the serialization of $\varphi$, $2^{|q|+1}\le B$, and $|x|\le B$, then for every natural
number $S<2^{|q|+1}$ there is a run, with all intermediate configurations satisfying the
register bound and the machine’s precondition, from the node state with depth $|q|$,
assignment stack $S$, and pointer $0$ to the return state carrying the truth value of
the quantified formula specification at depth $|q|$ and stack $S$, with pointer $0$.
\end{theorem}
